\documentclass[12pt, oneside]{book}

% --- PACKAGES ---
\usepackage{amsmath}       % Math equations
\usepackage{dirtree}
\usepackage{listings}
\usepackage{xcolor}        % Colored text
\usepackage{tcolorbox}     % Boxed text/equations
\usepackage{tikz}          % Drawing and image overlays
\usepackage{graphicx}      % Images
\usepackage[colorlinks=true, linkcolor=blue]{hyperref} % Interactive PDF links
\usepackage[T1]{fontenc}
\usepackage{amssymb}
\usepackage{amsfonts}
\usepackage{bbm}

\lstset{
    language=Python,
    basicstyle=\ttfamily\small,
    keywordstyle=\color{blue}\bfseries,
    commentstyle=\color{green!50!black}\itshape,
    stringstyle=\color{red},
    showstringspaces=false,
    numbers=left,
    numberstyle=\tiny\color{gray},
    frame=lines,
    backgroundcolor=\color{black!5},
    breaklines=true,
    mathescape=true, % Allows $...$ math inside comments to render properly
    morekeywords={fn, class, pass, return, if, elif, else, for, in} % Adds your custom keywords
}

\begin{document}

\pagenumbering{roman}
\begin{titlepage}
    \centering
    
    \vspace*{3cm} 
    
    {\Huge \textbf{FEM Basics}\par}
    
    \vspace*{3cm} 
    \vfill
    
    {\small Oliver Jackson\par}
    
    \vspace*{0.5cm} 

    {\small \today\par}
\end{titlepage}


\tableofcontents

\cleardoublepage
\pagenumbering{arabic}

\chapter{Introduction}

This document details the FEM Galerkin method. I included both the derivation of the method and am developing a 
software package to implement the method in Python.

\section{The weak form}

We are considering problems of the ellipcic type in one dimension. 

$$
-\frac{d}{d x}\left(c(x) \frac{d u}{d x}\right)=f(x) \quad \text { with solution } u.
$$

The space $c^{k}(\bar{\Omega})$ is defined to be the set of all real-valved functions 
a defined on $\bar{\Omega}$ (which is both $\Omega$ and $\partial \Omega$ ) with the property that 
$u$ and its derivatives up to the $k^{\text {th }}$ order are all continuous on $\bar{\Omega}$. 
A solution to $u$ is sought in the subspace

$$
C_{D}^{2}(\bar{\Omega})=\left\{u \in C^{2}(\bar{\Omega}): u=0 \text { on } \partial \Omega\right\} .
$$

Introduce the test function $v$ and the space of test functions to be $C_{D}^{2}(\bar{\Omega})$.

$$
-\nabla(c \cdot \nabla u)=f
$$

$$
\int-\nabla(c \cdot \nabla u) \cdot v=\int f \cdot v.
$$

This equation holds for some $u \in C_{D}^{2}(\bar{\Omega})$ denoted $u_{h}$ and for all 
$v \in C_{D}^{2}(\bar{\Omega})$ if and only if $u$ satisfies the original differential 
equation. The next step is to apply green's identity to the left hand side. A quick 
review of Green's identity. The weak form of a BVP is derived using greens 
identity, which is the multi-dimensional form of integration by parts.

$$
\nabla \cdot(v \nabla u)=\nabla v \cdot \nabla u+v \Delta u.
$$

Then to obtain Greens identity integrate both sides over $\Omega$ and apply the divergence theorem

$$
\begin{aligned}
& \Rightarrow \int_{\Omega} \nabla \cdot(v \nabla u)=\int_{\Omega} \nabla v \cdot \nabla u+\int_{\Omega} v \Delta u \\
& \Rightarrow \int_{\partial \Omega} v \nabla u \cdot n=\int_{\Omega} \nabla v \cdot \nabla u+\int_{\Omega} v \Delta u \\
& \Rightarrow-\int_{\Omega} v \Delta u=\int_{\Omega} \nabla v \cdot \nabla u-\int_{\partial \Omega} v \frac{\partial u}{\partial n}.
\end{aligned}
$$

we know $V=0$ on $\partial \Omega$ so

$$
-\int_{\Omega} v \Delta u=\int_{\Omega} \nabla V \cdot \nabla u.
$$

So we get

$$
\int_{\Omega}-\nabla(c \cdot \nabla u) \cdot v=\int_{\Omega} c \nabla u \cdot \nabla v
$$

So the weak form of our elliptic problem is given by Find $u \in C_{D}^{2}(\bar{\Omega})$ such 
that $\int_{\Omega} c \nabla u \cdot \nabla v=\int_{\Omega} fv$ for all $v \in C_{D}^{2}(\bar{\Omega})$.
We introduce notation $(u, v)=\int_{\Omega} u \cdot v$ and bilinear operator $a$,

$$
a(u, v)=(f, v) \quad \forall v \in C_{D}^{2}(\bar{\Omega})
$$

Where $a(u, v) \equiv \int_{\Omega} c \nabla u \cdot \nabla v$. 

\section{Weak form over a finite number of nodes}

Define nodes on the interval $x_{0}<x_{1}<\cdots<x_{N}<x_{N+1}$. The space between the nodes $N$ are $h=x_{n}-x_{n-1}$ 
where $n=0, \ldots, N+1$. As parameters to describe a function $v \in C_{D}^{2}(\bar{\Omega})$ we may choose the 
values $\eta_{n}=V\left(x_{n}\right)$ at the node points $n=0, \ldots, N+1$. To construct functions $v \in C_{D}^{2}(\bar{\Omega})$ 
we can define basis functions such that overlapping and neighboring hodes return a value one and disconnected nodes return zero.
Linear, quadratic, or any other valid basis functions can be used to formulate these functions over each interval. Here we 
define basis functions such that

$$
\phi_{j}\left(x_{i}\right)= \begin{cases}1 & \text { if } \quad i=j \\ 0 & \text { if } \quad i \neq j \quad i, j=1, \ldots, N.\end{cases}
$$

A function $v \in C_{D}^{2}(\bar{\Omega})$ then has the representation

$$
v(x)=\sum_{m=1}^{M} \eta_{m} \phi_{m}(x) \quad x \in[0,1].
$$

Then the FEM can now be formulated as: Find $u_{h} \in C_{D}^{2}(\bar{\Omega})$ such that

$$
a\left(u_{h}, v\right)=(f, v) \quad \forall v \in C_{D}^{2}(\bar{\Omega}).
$$

If $u_{h}$ satisfies the above then

\begin{equation}
a\left(u_{h}, \phi_{m}\right)=\left(f, \phi_{m}\right) \quad m=1, \ldots, M  \tag{*}
\end{equation}

Then by taking linear combinations we see $u_h$ satisfies $\xi$ since

$$
u_{h}(x)=\sum_{n=1}^{N} \xi_{n} \phi_{n}(x), \quad \xi_{n}=u_{h}\left(x_{n}\right)
$$

rewrite * like so

$$
\sum_{n=1}^{N} \xi_{n} a\left(\phi_{n}, \phi_{m}\right)=\left(f, \phi_{m}\right) \quad m=1, \ldots, M
$$

which is a linear system of equations with M equations and $N$ unknowus $\xi_{1}, \ldots, \xi_{N}$ which 
we are solving for. In matrix form the linear system of equations can be written as

$$
A \xi=b
$$

Where $A$ is the $M \times N$ matrix with elements $A_{m n}=a\left(\phi_{n}, \phi_{m}\right)$ 
called stiffness matrix. And where $\xi_n=\left(\xi_{1}, \ldots, \xi_{N}\right)$ and $b_m=\left(b_{1}, \ldots, b_{M}\right)$ 
with $b_{m}=\left(f, \phi_{m}\right)$ are both column vectors and $b$ is the load vector. Row $m$ is test and column $n$ 
is trail. 

\section{General algorithm for stiffness matrix construction}

Now we want to build a general algorithm to construct the stiffness matrix which decouples the geometric 
shape of the mesh from the functional degrees of freedom (the nodes). The algorithm will be general in the 
sense that any number of nodes $N$ can be chosen and the geometry, dimension, and basis type are separated through 
if/else statements in the code. Further calculus and Jacobian matrix is required if a general function for the 
construction of mesh element and nodes matrices is desired. Instead, the general framework advised by Prof He is to 
add additional definitions for linear VS quadratic basis functions or $P_b$ and $T_b$ information matrices into if/else 
statements. Inside the assembly code we will call these components (if/else statements) to construct the stiffness matrix.
\par The $P$ matrix is an information matrix consisting of the coordinates of all nodes. It is a $D \times N$ matrix 
where $D$ is dimension. In $1 D$ and $2 D$

$$
\begin{aligned}
P_{1 D} & =\left(x_{1}, \ldots, x_{N}\right) \\
P_{2 D} & =\binom{x_{1}, \ldots, x_{N}}{y_{1}, \ldots, y_{N}}
\end{aligned}
$$

The $T$ matrix is an information matrix containing the indices of the nodes of each element. 
Here the elements are the integral bounds over each interval. So the $T$ matrix holds the indices 
of the elements $E_{k}$ where $k$ denotes the number of elements $k=1, \ldots, N-1$. 
Define $N_b$ to be the number of basis functions required to construct one element. Then $T$ is a $N_b \times K$ matrix. 
\par For the linear basis function case there are two basis functions per element so $N_b=2$ and $K=N-1$. 
Then the $T$ matrix is given by

$$
T=\left(\begin{array}{lllll}
1 & 2 & \cdots & N-2 & N-1 \\
2 & 3 & \cdots & N-1 & N
\end{array}\right). 
$$

Only when we have linear basis functions we can define

\begin{center}
  \textit{contains $\Omega$ points $\mathbf{P=P_{b}}$ where $b$ denotes basis node value}. \\
  \textit{contains $\Omega$ elements $\mathbf{T=T_{b}}$ where $b$ denotes basis element value}.
\end{center}

When our domain is defined with quadratic basis functions instead of linear then the number of points needed per element
increases but NOT the number of nodes $N$ in our domain. Then $P_b$ will contain the midpoints in addition to the two endpoints 
to describe the quadratic curve. $T$ and $T_b$ then must also have an extra row so we know where this extra midpoint is physically 
located in terms of its neighbors. Thus $P$ and $T$ contain information about the geometry of the system which is represented in the 
element stiffness matrix. $T_{b}$ contains algebraic information used to map the element stiffness 
matrix to the global stiffness matrix. $P_b$ contains algebraic information about basis function construction. Then for 
quadratic basis $P \neq P_{b}$ and $T \neq T_{b}$. 
\par Now consider the interval of a specific element 
$E=\left[x_{n}, x_{n+1}\right]$. Then because our basis functions are defined for 
$i \neq j$ such that $\phi_{i}\left(x_{j}\right)=0$ and $\phi_{i}\left(x_{i}\right)=1$, the only basis functions we 
consider over this element $E$ are $\phi_{n}, \phi_{n+1}$. Note the following important property

$$
\left(\phi_{j}, \phi_{j+1}\right)=\left(\phi_{j+1}, \phi_{j}\right).
$$

Then we define local basis function $\psi$ as $\phi$ over the interval E.
Recall the linear basis functions:

$$
\phi_{n}(x)=\left\{\begin{array}{cl}
\frac{x-x_{n-1}}{h} & \text { for } x \in\left[x_{n-1}, x_{n}\right] \\
\frac{x_{n+1}-x}{h} & \text { for } x \in\left[x_{n}, x_{n+1}\right] \\
0 & \text { else }
\end{array}\right.
$$

Then define the local basis functions for an element $E=\left[x_{n}, x_{n+1}\right]$. On this element we have 
contributions from the positive edge of $\phi_{n+1}$ and negative from $\phi_{n}$. Call the negative slope $\psi_{1}$  
and the positive slope $\psi_{2}$ given by

$$
\begin{aligned}
& \psi_{1}=\frac{x_{n+1}-x}{h} \text { for } x \in\left[x_{n}, x_{n+1}\right] \\
& \psi_{2}=\frac{x-x_{n}}{h} \text { for } x \in\left[x_{n}, x_{n+1}\right]. 
\end{aligned}
$$

So on one element $E$ we have two local basis functions denoted $N_b=2$ for linear global basis functions 
and the total number of global basis functions is $N$. We need two points $\left[x_{n}, x_{n+1}\right]$ to 
compute $\psi_{1}, \psi_{2}$ that are extracted directly from the $P_{b}$ matrix using the indices from the Tb matrix:

$$
\begin{aligned}
& x_{n}^{(k)}=P_{b}\left(1, T_{b}(1, k)\right) \\
& x_{n+1}^{(k)}=P_{b}\left(1, T_{b}(2, k)\right)
\end{aligned}
$$

There are four non-zero local integrals on $E_{k}=\left[x_{n}, x_{n+1}\right]$

$$
\begin{aligned}
& a_{\beta \alpha}^{k}=a^{k}\left(\psi_{\beta}, \psi_{\alpha}\right)  \\
& a_{\beta \alpha}^{k}=\int_{x_{n}^{k}}^{x_{n+1}^{k}} c\left(\psi_{\beta}^{\prime} \psi_{\alpha}^{\prime}\right) d x \text { for } \beta, \alpha \in\{1,2\}
\end{aligned}
$$

This means the local integral should be assembled to the global matrix entry $a_{m n}$. An element is given 
by $E_{k}=\left[x_{n}^{k}, x_{n+1}^{k}\right]$ for $k=1, \ldots, N-1$. We compute an $N_{b} \times N_{b}$ 
element stiffness matrix $a^{k}$ with indices $\alpha, \beta \in\{1,2\}$ for the linear case. Because $\beta$ 
is the test function index it indicates the rows of $a^{k}$ and $\alpha$ the columns.

$$
a_{\beta \alpha}^{k}=\left[\begin{array}{cc}
a_{11}^{k} & a_{12}^{k} \\
a_{21}^{k} & a_{22}^{k}
\end{array}\right]=\left[\begin{array}{ll}
\int\left(\psi_{1}^{\prime} \psi_{1}^{\prime}\right. & \int c \psi_{1}^{\prime} \psi_{2}^{\prime} \\
\int c \psi_{2}^{\prime} \psi_{1}^{\prime} & \int c \psi_{2}^{\prime} \psi_{2}^{\prime}
\end{array}\right]
$$

To map the local element stiffness matrix $a_{\beta \alpha}^{k}$ into the global $M \times N$ matrix $A$ we 
use the topology matrix $T_{b}$ to route the indices.

$$
\begin{aligned}
& \text { Global row } m=T_{b}(\beta, k) \\
& \text { Global column } n=T_{b}(\alpha, k)
\end{aligned}
$$

I write the global entry of $A$ as $A_{m n}$ and the local element stiffness matrix as $a_{\beta \alpha}^{k}$. I 
include $\delta_{m, T_{b}(\beta, k)}$ as a way to search for the global row index and $\delta_{n}, T_{b}(\alpha, k)$ 
as a way to search for the global column index.

$$
A_{m n}=\sum_{n=1}^{N-1} \sum_{\beta=1}^{N_{b}} \sum_{\alpha=1}^{N_{b}} a_{\beta \alpha}^{k} \cdot \delta_{m, T_{b}(\beta, k)} \cdot \delta_{n, T_{b}(\alpha, k)}
$$

\chapter{Programming}
\section{File structure}
To program this we can utilize the following file structure: 

\vspace{1cm}
\dirtree{%
.1 FEMelliptic/.
.2 src/.
.3 main.py.
.4 build\_stiffness\_matrices.py.
.4 gaussian\_quadrature.py.
.4 boundary\_conditions.py.
.4 error.py.
.4 basis.py.
.4 information\_matrices.py.
.4 functions.py.
.3 documentation/.
.4 FEM\_basics.tex.
}
\vspace{1cm}

The assembly code contains the general algorithm for stiffness matrix construction. The components contains the 
if/else statements to define the geometry and basis functions. The main.py file is the driver code that 
calls the assembly code and the components to construct the stiffness matrix.

\section{main.py}

The build\_stiffness\_matrices.py file contains the general algorithms provided for $A_{mn}$ and $b_m$. Here, instead of dirac delta functions we can use for loops to search for the row and column indices of the stiffness matrix. 

In main.py we will define the following parameters:

\begin{lstlisting}
N = int # number of internal nodes in domain $N$ which corresponds to the number of unknowns \\
K = N+1 # number of elements \\
N\_b = int # number of basis functions per element \\
D = int # dimension of the problem \\
gauss\_point\_number = int # number of gauss points for integration \\

x = linspace(0, 1, N+2) # domain \\

stiffness\_matrix = build\_stiffness\_matrix(N, N\_b, K, P\_b, T\_b, c\_func, basis\_type, gauss\_point\_number) \\
load\_matrix = build\_load\_matrix(N, P\_b, T\_b, f\_func, basis\_type, gauss\_point\_number) \\

left\_boundary = 0 # homogeneous Dirichlet boundary condition \\
right\_boundary = 0
stiffness\_matrix, load\_matrix = apply\_dirichlet\_linear(stiffness\_matrix, load\_matrix, left\_boundary, right\_boundary) \\

analytical\_solution = functions.solution\_func(x) \\
FEM\_solution = np.linalg.solve(stiffness\_matrix, load\_matrix) \\
error = error(analytical\_solution.flatten(), FEM\_solution.flatten()) \\
\end{lstlisting}

\section{build\_stiffness\_matrices.py}

build\_stiffness\_matrices.py contains the following functions:

\begin{lstlisting}
fn build\_local\_stiffness\_matrix(c\_func, basis\_type, left\_node, right\_node, gauss\_point\_number, N\_b) \\
    local\_stiffness\_matrix = zeros(N\_b, N\_b) \\
    for beta in range(1, N\_b + 1): \\
        for alpha in range(1, N\_b + 1): \\
            fn integrand(x) \\
                d\_psi\_beta = basis\_type.derivative(beta, left\_node, right\_node, x) \\
                d\_psi\_alpha = basis\_type.derivative(alpha, left\_node, right\_node, x) \\
                return c\_func(x) * d\_psi\_beta * d\_psi\_alpha \\
            local\_stiffness\_matrix[beta - 1, alpha - 1] = gauss\_integrate(
                integrand, 
                left\_node, 
                right\_node, 
                gauss\_point\_number) \\
    return local\_stiffness\_matrix # returns the local stiffness matrix for one element (2x2 for linear) \\


fn build\_stiffness\_matrix(N, N\_b, K, P\_b, T\_b, c\_func, basis\_type, gauss\_point\_number) \\
    stiffness\_matrix = sparse(N+2, N+2) \\
    for k in range(K): \\
        left\_node = P\_b[0, T\_b[0, k]] \\
        right\_node = P\_b[0, T\_b[1, k]] \\
        local\_stiffness\_matrix = build\_local\_stiffness\_matrix(
            c\_func,
            basis\_type, # basis.linear or basis.quadratic
            left\_node,
            right\_node,
            gauss\_point\_number)\\
        for beta in range(1, N\_b + 1): \\
            for alpha in range(1, N\_b + 1): \\
                stiffness\_matrix[T\_b[beta - 1, k], T\_b[alpha - 1, k]] += local\_stiffness\_matrix[beta - 1, alpha - 1] \\
    return stiffness\_matrix # returns the global stiffness matrix (NxN) \\

fn build\_load\_matrix(N, N\_b, K, P\_b, T\_b, f\_func, basis\_type, gauss\_point\_number) \\
    load\_matrix = zeros(N+2, 1) \\
    for k in range(K): \\
        left\_node = P\_b[0, T\_b[0, k]] \\
        right\_node = P\_b[0, T\_b[1, k]] \\
        for beta in range(1, N\_b + 1): \\
            fn integrand(x) \\
                psi\_beta = basis\_type.\_\_call\_\_(beta, left\_node, right\_node, x)
                return f\_func(x) * psi\_beta \\
            load\_matrix[T\_b[beta - 1, k]] += gauss\_integrate(
                integrand, 
                left\_node, 
                right\_node, 
                gauss\_point\_number
            ) \\
    return load\_matrix # returns the global load matrix (Nx1) \\
\end{lstlisting}

\section{basis.py}
            
File for basis functions basis.py contains the following functions:

\begin{lstlisting}
class linear
    fn derivative(beta, left\_node, right\_node) \\
        h = right\_node - left\_node \\
        if beta == 1: \\
            return -1/h \\
        elif beta == 2: \\
            return 1/h \\
        else: \\
            raise ValueError("Invalid beta index for linear basis function") \\
    fn \_\_call\_\_(beta, left\_node, right\_node, x) \\
        h = right\_node - left\_node \\
        if beta == 1: \\
            return (right\_node - x) / h \\
        elif beta == 2: \\
            return (x - left\_node) / h \\
class quadratic
    fn derivative(beta, left\_node, right\_node, x) \\
        h = right\_node - left\_node \\
        if beta == 1: \\
            return (1/h) * (4 * (x - left\_node) - 3) \\
        elif beta == 2: \\
            return (1/h) * (4 * (x - left\_node) - 1) \\
        elif beta == 3: \\
            return (1/h) * (-8 * (x - left\_node) + 4) \\
        else: \\
            raise ValueError("Invalid beta index for quadratic basis function") \\
    fn \_\_call\_\_(beta, left\_node, right\_node, x) \\
        h = right\_node - left\_node \\
        if beta == 1: \\
            return 2 * ((x - left\_node)/h)**2 - 3 * ((x - left\_node)/h) + 1
        if beta == 2: \\
            return 2 * ((x - left\_node)/h)**2 - ((x - left\_node)/h)
        if beta == 3: \\
            return -4 * ((x - left\_node)/h)**2 + ((x - left\_node)/h)
\end{lstlisting}

\section{gaussian\_quadrature.py}

File for guassian quadrature gaussian\_quadrature.py contains the following functions:

\begin{lstlisting}
fn generate\_gauss\_reference\_1D(gauss\_point\_number)
    if gauss\_point\_number == 2:
        gauss\_coefficient\_reference\_1D = [1.0, 1.0] \\
        gauss\_point\_reference\_1D = [-1.0/sqrt(3), 1.0/sqrt(3)]
    elif gauss\_point\_number == 4:
        gauss\_coefficient\_reference\_1D = [0.3478548451, 
        0.3478548451, 0.6521451549, 0.6521451549] \\
        gauss\_point\_reference\_1D = [0.8611363116, 
        -0.8611363116, 0.3399810436, -0.3399810436] \\
    elif gauss\_point\_number == 8:
        gauss\_coefficient\_reference\_1D = [0.1012285363, 
        0.1012285363, 0.2223810345, 0.2223810345, 0.3137066459, 
        0.3137066459, 0.3626837834, 0.3626837834] \\
        gauss\_point\_reference\_1D = [0.9602898565, 
        -0.9602898565, 0.7966664774, -0.7966664774, 0.5255324099, 
        -0.5255324099, 0.1834346425, -0.1834346425] \\
    else:
        raise ValueError("Unsupported number of Gauss points. Choose 2, 4, or 8.")
    return gauss\_coefficient\_reference\_1D, gauss\_point\_reference\_1D

fn generate\_gauss\_local\_1D(gauss\_coefficient\_reference\_1D, gauss\_point\_reference\_1D, lower\_bound, upper\_bound)
    gauss\_coefficient\_local\_1D = []
    gauss\_point\_local\_1D = []

    for i in range(len(gauss\_coefficient\_reference\_1D))
        local\_coeff = (upper\_bound - lower\_bound) * gauss\_coefficient\_reference\_1D[i] / 2.0
        local\_point = (upper\_bound - lower\_bound) * gauss\_point\_reference\_1D[i] / 2.0 + (upper\_bound + lower\_bound) / 2.0
        gauss\_coefficient\_local\_1D.append(local\_coeff)
        gauss\_point\_local\_1D.append(local\_point)
    return gauss\_coefficient\_local\_1D, gauss\_point\_local\_1D

fn gauss\_integrate(integrand\_func, left\_node, right\_node, gauss\_point\_number)
    ref\_coeffs, ref\_points = generate\_gauss\_reference\_1D(gauss\_point\_number)
    local\_coeffs, local\_points = generate\_gauss\_local\_1D(ref\_coeffs, ref\_points, left\_node, right\_node)
    integral = 0.0
    for i in range(gauss\_point\_number)
        integral += local\_coeffs[i] * integrand\_func(local\_points[i])
    return integral
\end{lstlisting}

\section{information\_matrices.py}

File for generating the P, T, P\_b, and T\_b information matrices 
for linear and quadratic 1D. The information matrices are stored in information\_matrices.py and contain the following functions:

\begin{lstlisting}
class linear 
    fn P(N)
        return np.array([np.linspace(0, 1, N+2)]) # returns the P matrix for 1D linear basis functions
    fn T(N)
        K = N + 1
        return np.array([
            [k for k in range(K)], 
            [k + 1 for k in range(K)]
        ]) # returns the T matrix for 1D linear basis functions
    fn P\_b\_linear(N)
        return linear.P(N) 
    fn T\_b\_linear(N)
        return linear.T(N)

class quadratic
    # Implement similar functions for quadratic basis functions
\end{lstlisting}

\section{functions.py}

functions.py contains the c\_funct and f\_func functions:

\begin{lstlisting}
class c\_func
    fn hw\_1(x)
        return exp(x) 
class f\_func
    fn hw\_1(x)
        return -exp(x) * (cos(x) - 2 * sin(x) - x * cos(x) - x * sin(x))
class solution\_func(x)
    return x * cos(x) 
\end{lstlisting}

\section{boundary\_conditions.py and error.py}

The boundary\_conditions.py file contains the following functions:

\begin{lstlisting}
class dirichlet
    fn apply\_dirichlet\_linear(stiffness\_matrix, load\_matrix, left\_boundary, right\_boundary)
        stiffness\_matrix[0, :] = 0
        stiffness\_matrix[0, 0] = 1
        load\_matrix[0] = left\_boundary

        stiffness\_matrix[-1, :] = 0
        stiffness\_matrix[-1, -1] = 1
        load\_matrix[-1] = right\_boundary

        return stiffness\_matrix, load\_matrix

class robin
    fn apply\_robin\_linear(stiffness\_matrix, load\_matrix)
        pass # Implement the Robin boundary condition logic here
\end{lstlisting}

The error.py file contains the following functions:

\begin{lstlisting}
fn error(analytical\_solution, FEM\_solution):
    abs\_max\_difference = np.max(np.abs(analytical\_solution - FEM\_solution))
    return abs\_max\_difference
\end{lstlisting}

       





NOTATION KEY

- $N=$ number of internal nodes in domain $N$ which corresponds to the \# of cukhowns \\
- $k=G$ lobal element index $(k=1, \ldots, N-1) k$ \\
- $m=$ Global test basis index $(1, \ldots, M)$ maps rows of $A M$ \\
- $n=$ Global trail basis index $(1, \ldots, N)$ maps to columns of $A n$ \\
- $\beta=$ local test basis index $\beta \in\{1,2\}$ beta \\
- $\alpha=$ local trail basis index $\alpha \in\{1,2\}$ alpha \\
- $\xi_{n}=$ global trial coefficients $x_{i}$ \\
- $\eta_{m}=$ global test coefficients eta \\
- $a_{\beta \alpha}^{k}=$ local stiffness matrix local-a \\
- domain $=X=x_{0}, \ldots, x_{N+1} \quad x$ \\
- $T, T_{b}=$ topology information matrix $T$-matrix, $T$ - $b$-matrix \\
- $P, P_{b}=$ index information matrix $P$-matrix, $P$ - $b$-matrix \\

\end{document}

\section{main.py}

The build\_stiffness\_matrices.py file contains the general algorithms provided for $A_{mn}$ and $b_m$. Here, instead of 
dirac delta functions we can use for loops to search for the row and column indices of the stiffness matrix. 

In main.py we will define the following parameters:

N = int # number of internal nodes in domain $N$ which corresponds to the number of unknowns \\
K = N+1 # number of elements \\
N\_b = int # number of basis functions per element \\
D = int # dimension of the problem \\
gauss\_point\_number = int # number of gauss points for integration \\

x = linspace(0, 1, N+2) # domain \\

stiffness\_matrix = build\_stiffness\_matrix(N, N\_b, K, P\_b, T\_b, c\_func, basis\_type, gauss\_point\_number) \\
load\_matrix = build\_load\_matrix(N, P\_b, T\_b, f\_func, basis\_type, gauss\_point\_number) \\

left\_boundary = 0 # homogeneous Dirichlet boundary condition \\
right\_boundary = 0
stiffness\_matrix, load\_matrix = apply\_dirichlet\_linear(stiffness\_matrix, load\_matrix, left\_boundary, right\_boundary) \\

analytical\_solution = functions.solution\_func(x) \\
FEM\_solution = np.linalg.solve(stiffness\_matrix, load\_matrix) \\
error = error(analytical\_solution.flatten(), FEM\_solution.flatten()) \\

\section{build\_stiffness\_matrices.py}

build\_stiffness\_matrices.py contains the following functions:

fn build\_local\_stiffness\_matrix(c\_func, basis\_type, left\_node, right\_node, gauss\_point\_number, N\_b) \\
    local\_stiffness\_matrix = zeros(N\_b, N\_b) \\
    for beta in range(1, N\_b + 1): \\
        for alpha in range(1, N\_b + 1): \\
            fn integrand(x) \\
                d\_psi\_beta = basis\_type.derivative(beta, left\_node, right\_node, x) \\
                d\_psi\_alpha = basis\_type.derivative(alpha, left\_node, right\_node, x) \\
                return c\_func(x) * d\_psi\_beta * d\_psi\_alpha \\
            local\_stiffness\_matrix[beta - 1, alpha - 1] = gauss\_integrate(
                integrand, 
                left\_node, 
                right\_node, 
                gauss\_point\_number) \\
    return local\_stiffness\_matrix # returns the local stiffness matrix for one element (2x2 for linear) \\


fn build\_stiffness\_matrix(N, N\_b, K, P\_b, T\_b, c\_func, basis\_type, gauss\_point\_number) \\
    stiffness\_matrix = sparse(N+2, N+2) \\
    for k in range(K): \\
        left\_node = P\_b[0, T\_b[0, k]] \\
        right\_node = P\_b[0, T\_b[1, k]] \\
        local\_stiffness\_matrix = build\_local\_stiffness\_matrix(
            c\_func,
            basis\_type, # basis.linear or basis.quadratic
            left\_node,
            right\_node,
            gauss\_point\_number)\\
        for beta in range(1, N\_b + 1): \\
            for alpha in range(1, N\_b + 1): \\
                stiffness\_matrix[T\_b[beta - 1, k], T\_b[alpha - 1, k]] += local\_stiffness\_matrix[beta - 1, alpha - 1] \\
    return stiffness\_matrix # returns the global stiffness matrix (NxN) \\

fn build\_load\_matrix(N, N\_b, K, P\_b, T\_b, f\_func, basis\_type, gauss\_point\_number) \\
    load\_matrix = zeros(N+2, 1) \\
    for k in range(K): \\
        left\_node = P\_b[0, T\_b[0, k]] \\
        right\_node = P\_b[0, T\_b[1, k]] \\
        for beta in range(1, N\_b + 1): \\
            fn integrand(x) \\
                psi\_beta = basis\_type.\_\_call\_\_(beta, left\_node, right\_node, x)
                return f\_func(x) * psi\_beta \\
            load\_matrix[T\_b[beta - 1, k]] += gauss\_integrate(
                integrand, 
                left\_node, 
                right\_node, 
                gauss\_point\_number
            ) \\
    return load\_matrix # returns the global load matrix (Nx1) \\


\section{basis.py}
            
File for basis functions basis.py contains the following functions:

class linear
    fn derivative(beta, left\_node, right\_node) \\
        h = right\_node - left\_node \\
        if beta == 1: \\
            return -1/h \\
        elif beta == 2: \\
            return 1/h \\
        else: \\
            raise ValueError("Invalid beta index for linear basis function") \\
    fn \_\_call\_\_(beta, left\_node, right\_node, x) \\
        h = right\_node - left\_node \\
        if beta == 1: \\
            return (right\_node - x) / h \\
        elif beta == 2: \\
            return (x - left\_node) / h \\
class quadratic
    fn derivative(beta, left\_node, right\_node, x) \\
        h = right\_node - left\_node \\
        if beta == 1: \\
            return (1/h) * (4 * (x - left\_node) - 3) \\
        elif beta == 2: \\
            return (1/h) * (4 * (x - left\_node) - 1) \\
        elif beta == 3: \\
            return (1/h) * (-8 * (x - left\_node) + 4) \\
        else: \\
            raise ValueError("Invalid beta index for quadratic basis function") \\
    fn \_\_call\_\_(beta, left\_node, right\_node, x) \\
        h = right\_node - left\_node \\
        if beta == 1: \\
            return 2 * ((x - left\_node)/h)**2 - 3 * ((x - left\_node)/h) + 1
        if beta == 2: \\
            return 2 * ((x - left\_node)/h)**2 - ((x - left\_node)/h)
        if beta == 3: \\
            return -4 * ((x - left\_node)/h)**2 + ((x - left\_node)/h)

\section{gaussian\_quadrature.py}

File for guassian quadrature gaussian\_quadrature.py contains the following functions:

fn generate\_gauss\_reference\_1D(gauss\_point\_number)
    if gauss\_point\_number == 2:
        gauss\_coefficient\_reference\_1D = [1.0, 1.0] \\
        gauss\_point\_reference\_1D = [-1.0/sqrt(3), 1.0/sqrt(3)]
    elif gauss\_point\_number == 4:
        gauss\_coefficient\_reference\_1D = [0.3478548451, 
        0.3478548451, 0.6521451549, 0.6521451549] \\
        gauss\_point\_reference\_1D = [0.8611363116, 
        -0.8611363116, 0.3399810436, -0.3399810436] \\
    elif gauss\_point\_number == 8:
        gauss\_coefficient\_reference\_1D = [0.1012285363, 
        0.1012285363, 0.2223810345, 0.2223810345, 0.3137066459, 
        0.3137066459, 0.3626837834, 0.3626837834] \\
        gauss\_point\_reference\_1D = [0.9602898565, 
        -0.9602898565, 0.7966664774, -0.7966664774, 0.5255324099, 
        -0.5255324099, 0.1834346425, -0.1834346425] \\
    else:
        raise ValueError("Unsupported number of Gauss points. Choose 2, 4, or 8.")
    return gauss\_coefficient\_reference\_1D, gauss\_point\_reference\_1D

fn generate\_gauss\_local\_1D(gauss\_coefficient\_reference\_1D, gauss\_point\_reference\_1D, lower\_bound, upper\_bound)
    gauss\_coefficient\_local\_1D = []
    gauss\_point\_local\_1D = []

    for i in range(len(gauss\_coefficient\_reference\_1D))
        local\_coeff = (upper\_bound - lower\_bound) * gauss\_coefficient\_reference\_1D[i] / 2.0
        local\_point = (upper\_bound - lower\_bound) * gauss\_point\_reference\_1D[i] / 2.0 + (upper\_bound + lower\_bound) / 2.0
        gauss\_coefficient\_local\_1D.append(local\_coeff)
        gauss\_point\_local\_1D.append(local\_point)
    return gauss\_coefficient\_local\_1D, gauss\_point\_local\_1D

fn gauss\_integrate(integrand\_func, left\_node, right\_node, gauss\_point\_number)
    ref\_coeffs, ref\_points = generate\_gauss\_reference\_1D(gauss\_point\_number)
    local\_coeffs, local\_points = generate\_gauss\_local\_1D(ref\_coeffs, ref\_points, left\_node, right\_node)
    integral = 0.0
    for i in range(gauss\_point\_number)
        integral += local\_coeffs[i] * integrand\_func(local\_points[i])
    return integral

\section{information\_matrices.py}

File for generating the P, T, P\_b, and T\_b information matrices 
for linear and quadratic 1D. The information matrices are stored in information\_matrices.py and contain the following functions:

class linear 
    fn P(N)
        return np.array([np.linspace(0, 1, N+2)]) # returns the P matrix for 1D linear basis functions
    fn T(N)
        K = N + 1
        return np.array([
            [k for k in range(K)], 
            [k + 1 for k in range(K)]
        ]) # returns the T matrix for 1D linear basis functions
    fn P\_b\_linear(N)
        return linear.P(N) 
    fn T\_b\_linear(N)
        return linear.T(N)

class quadratic
    # Implement similar functions for quadratic basis functions

\section{functions.py}

functions.py contains the c\_funct and f\_func functions:

class c\_func
    fn hw\_1(x)
        return exp(x) 
class f\_func
    fn hw\_1(x)
        return -exp(x) * (cos(x) - 2 * sin(x) - x * cos(x) - x * sin(x))
class solution\_func(x)
    return x * cos(x) 

\section{boundary\_conditions.py and error.py}

The boundary\_conditions.py file contains the following functions:

class dirichlet
    fn apply\_dirichlet\_linear(stiffness\_matrix, load\_matrix, left\_boundary, right\_boundary)
        stiffness\_matrix[0, :] = 0
        stiffness\_matrix[0, 0] = 1
        load\_matrix[0] = left\_boundary

        stiffness\_matrix[-1, :] = 0
        stiffness\_matrix[-1, -1] = 1
        load\_matrix[-1] = right\_boundary

        return stiffness\_matrix, load\_matrix

class robin
    fn apply\_robin\_linear(stiffness\_matrix, load\_matrix)
        pass # Implement the Robin boundary condition logic here

The error.py file contains the following functions:

fn error(analytical\_solution, FEM\_solution):
    abs\_max\_difference = np.max(np.abs(analytical\_solution - FEM\_solution))
    return abs\_max\_difference

